%=============================================================================!
% 16.12  WHAT THE TRAJECTORY LOOKS LIKE                                       !
%=============================================================================!
\section{What the trajectory looks like}
\label{sec:ob-trajectory}

The observables of this chapter are computed after the run, from saved
frames, and most of their faults are made then: in the frames chosen, in
how they are prepared, or in where a curve is read. A healthy set of
observables checks itself through the identities derived above. For the
liquid argon of this chapter, $g(r)$ with the pair potential gives the
run's energy and pressure, to 0.6 and 0.2 of their errors (\cref{sec:ob-rdf}); $S(G)$ from the positions agrees with
the transform of $g(r)$ where the cut at half the cell does not matter
and tends to $\rho\kB T\kappa_T$ at small $G$ (\cref{sec:ob-sq}); the slope
of the MSD and the integral of $C_{vv}$ give the same $D$
(\cref{sec:ob-vacf}); the density of states has the value $12mD/\kB T$
at zero frequency (\cref{sec:ob-vdos}); and where
the free energy is known exactly, the histogram reproduces it within its
errors (\cref{sec:ob-fes}). Each of these fails, in a recognisable way,
when something has gone wrong. \Cref{fig:ob-faults} and
\cref{tab:ob-faults} show six faults on liquid argon and gather those of
earlier sections.

\textit{Wrapped positions.} Positions put back into the cell, as most
trajectory files store them, give an MSD that jumps whenever an atom
crosses a face. From the run of \cref{sec:ob-msd} wrapped every
\SI{100}{\femto\second}, it is \SI{167.8}{\angstrom\squared} at
\SI{20}{\pico\second} against the true 45.2, and 222.0 and
\SI{258.0}{\angstrom\squared} at 50 and \SI{100}{\pico\second}, levelling
off towards \SI{270}{\angstrom\squared}, $L^2/2$, the mean squared distance
between two points placed at random in the cell (\cref{fig:ob-faults}a). The check is
the largest move between frames, which must be a small fraction of the
cell (\texttt{diagnostics.largest\_jump}, \cref{sec:cs-trajectory}), and
the cure is unwrapping (\texttt{diagnostics.unwrap}).

\textit{A drift of the centre of mass.} Starting velocities drawn without
removing their total momentum give the centre of mass a speed of about
$\sqrt{3\kB T/Nm}$, \SI{1.80e-4}{\angstrom\per\femto\second} here, which
it keeps for ever at fixed energy. Added to the same paths, it makes
the MSD curve upwards and the fit over 2 to \SI{20}{\pico\second} give $D$
31\% too high (\cref{fig:ob-faults}b); removing the centre of mass, as
\texttt{msd} does by default, restores \num{3.764e-4}. The check is the
path of the centre of mass (\texttt{diagnostics.centre\_of\_mass\_path}).

\textit{A fit in the ballistic part.} A window that ends before the motion
has forgotten its velocity fits a straight line to a curve: over 0.02 to
\SI{0.2}{\pico\second} it gives $D$ a third too low, while windows starting
at \SI{0.05}{\pico\second} or later and ending ten times later agree with
the 2 to \SI{20}{\pico\second} value within 6\% (\cref{fig:ob-faults}c). The
check is the slope of $\ln\text{MSD}$ against $\ln t$, which must be 1
over the window.

\textit{$g(r)$ past half the cell.} Separations reduced to their nearest
image cannot exceed half the cell's width along any axis, so beyond
$L/2$ the shells are partly empty. Computed out to \SI{16}{\angstrom} in
the 256-atom cell of side \SI{23.26}{\angstrom}, $g$ falls to 0.50 at
\SI{14}{\angstrom}, where the 2048-atom liquid shows 1.010
(\cref{fig:ob-faults}d). \texttt{rdf} refuses a range beyond half the
smallest width, as the cell list does.

\textit{Transport under a local thermostat.} Langevin friction of
\SI{10}{\per\pico\second} damps every velocity: $C_{vv}$ falls to half its
first value in \SI{60}{\femto\second} instead of 160, and to zero in 170
instead of 320, and its integral gives $D = \SI{0.682e-4}{\angstrom\squared\per
\femto\second}$, 82\% below the value at fixed energy
(\cref{fig:ob-faults}e). Transport coefficients and spectra come from runs
at fixed energy or under a global thermostat such as CSVR, which
Chapter~15 found leaves $D$ unchanged within its error; under friction
they belong to the friction (\cref{sec:th-langevin}).

\textit{A Green-Kubo integral read too late.} The integral of a
correlation function settles when the correlation has died and then
wanders as noise accumulates. For one \SI{1}{\nano\second} run of 256
atoms the shear viscosity read at \SI{2}{\pico\second}, on the plateau, is
\SI{2.18e-4}{\pascal\second}, 18\% above the \num{1.853e-4} from all
\SI{14}{\nano\second}, the spread of a single run; read later, at 10, 20
and \SI{40}{\pico\second}, it is 1.91, 1.78 and \num{2.01e-4}
(\cref{fig:ob-faults}f), wandering by as much again. The check is to read
the plateau where the correlation has died, here 1 to
\SI{3}{\pico\second}, and to take the error from independent runs.

\begin{figure}[htb]
\centering
\includegraphics{ch16_observables/faults}
\caption{Six faults on liquid argon, each against the correct curve
(teal). (a) The MSD from positions wrapped into the cell. (b) The MSD of
paths whose centre of mass drifts at $\sqrt{3\kB T/Nm}$, not removed. (c)
$D$ fitted to windows from $t_1$ to $10t_1$, against $t_1$, with the 2 to
\SI{20}{\pico\second} value (dotted). (d) $g(r)$ of the 256-atom cell taken
past half its side (dotted) with the nearest image, against the
2048-atom $g(r)$. (e) $C_{vv}$ under Langevin friction of
\SI{10}{\per\pico\second}. (f) The Green-Kubo integral of the shear stress
from one \SI{1}{\nano\second} run, read out to \SI{40}{\pico\second},
against the value from \SI{14}{\nano\second} (dotted).}
\label{fig:ob-faults}
\end{figure}

\begin{table}[htb]
\centering
\small
\caption{Faults of observables, the check that finds each, and its size
in this chapter's runs.}
\label{tab:ob-faults}
\begin{tabular}{@{}>{\raggedright\arraybackslash}p{0.29\linewidth}
   >{\raggedright\arraybackslash}p{0.36\linewidth}
   >{\raggedright\arraybackslash}p{0.29\linewidth}@{}}
\toprule
fault & check & measured\\
\midrule
wrapped positions in an MSD & largest move between frames; unwrap &
\SI{168}{\angstrom\squared} at \SI{20}{\pico\second}, not 45\\
centre of mass drifting & path of the centre of mass; remove it & $D$
31\% high\\
fit in the ballistic part & slope of $\ln$MSD against $\ln t$ is 1 & $D$
32\% low\\
$g(r)$ past half the cell & range at most half the smallest width & $g =
0.50$ at \SI{14}{\angstrom}, not 1.01\\
transport under friction & run at fixed energy or with CSVR & $D$ 82\%
low at \SI{10}{\per\pico\second}\\
Green-Kubo read too late & read the plateau; errors from runs &
$\eta_{\mathrm s}$ from 1.78 to \num{2.18e-4}\\
bonds from one cutoff & band of hysteresis; least lifetime & 3.2 times
the events (\cref{sec:ob-bonds})\\
frames too sparse for a spectrum & stride below $1/2\nu_{\max}$; same
stride for both & overlap 0.795 at \SI{30}{\femto\second}
(\cref{sec:ob-vdos})\\
errors of $F$ from short blocks & errors over independent copies &
squared misses over squared errors 3.26, not 1 (\cref{sec:ob-fes})\\
energy distributions from two ensembles & same thermostat for both &
overlap 0.882, not 0.980 (\cref{sec:ob-fes})\\
\bottomrule
\end{tabular}
\end{table}

\begin{keyidea}
Observables check one another: $g(r)$ gives the energy and the pressure,
the MSD and the velocity autocorrelation give the same $D$, and the
density of states holds $D$ at zero frequency. Their common faults
are wrapped positions, an unremoved drift, a fit window in the ballistic
part, a range past half the cell, transport under a local thermostat, an
integral read too late, frames too sparse, and error bars from blocks
shorter than the slowest motion.
\end{keyidea}

\notebook{16}{break each observable and find the check that catches it}

\begin{selfcheck}
An MSD from a long run levels off at a value close to $L^2/2$. What is
the most likely fault?
\end{selfcheck}
